% Source: https://artofproblemsolving.com/community/c3h1503216_contest_8th_grade
% Stable identity: aops-topic-1503216-q01
Let the centres be $O_1$, $O_2$ and the intersection points $P$, $Q$. Then $O_1O_2=O_1P=O_2P=6$, so triangles $O_1O_2P$ and $O_1O_2Q$ are equilateral and $\angle PO_1Q=120^\circ$:

\begin{asy}
usepackage("amsmath");
size(6cm);
defaultpen(fontsize(10pt));
real r=6; pair O1=(0,0),O2=(6,0),P=(3,3*sqrt(3)),Q=(3,-3*sqrt(3));
fill(arc(O1,r,-60,60)--arc(O2,r,120,240)--cycle,paleblue);
draw(circle(O1,r),linewidth(0.9)); draw(circle(O2,r),linewidth(0.9));
draw(O1--P--O2--Q--cycle,red); draw(P--Q,dashed);
dot(O1);dot(O2);dot(P);dot(Q);
label("$O_1$",O1,W);label("$O_2$",O2,E);label("$P$",P,N);label("$Q$",Q,S);
label("$6$",(O1+P)/2,dir(-30),fontsize(9pt));
draw(arc(O1,1.2,-60,60),red); label("$\angle PO_1Q=120^\circ$",(3,-6.6),S,red);
label("\mbox{}",(min(currentpicture,user=true)+max(currentpicture,user=true))/2,black); // keeps the text after the figure black when LaTeX typesets the labels inline
\end{asy}

The common region (shaded) is cut by the chord $PQ$ into two congruent circular segments. The segment on the side of $O_2$ is the sector $O_1PQ$ minus the triangle $O_1PQ$:
\[
\dfrac{120}{360}\cdot\pi\cdot6^2-\dfrac{1}{2}\cdot6\cdot6\cdot\sin120^\circ=12\pi-9\sqrt3 .
\]
So the common area is
\[
2\left(12\pi-9\sqrt3\right)=24\pi-18\sqrt3 .
\]
Here $m=24$ and $n=18$; they are uniquely determined, because $\dfrac{\pi}{\sqrt3}$ is irrational. Hence
\[
\dfrac{m+n}{2}=\dfrac{24+18}{2}=\boxed{21}.
\]
